% Transcription — fonds Alexandre Grothendieck, shelfmark 129, batch 4 (pages 61-80).
% Established from the facsimile archives/batches/129/batch-04.pdf, prepared
% from a copy of 129.pdf fetched through the project's relay
% (https://grothendieck-relay.onrender.com/source/129.pdf; PDF page k =
% archivists' page 60+k, no cover sheet), per the skill transcribe-grothendieck.
% The folder is 114 pages in six batches, transcribed in parallel by six
% passes; this is the fourth (pages 61-80).
%
% Pass: Opus 5.5 (claude-opus-5-5), 2026-09-26 — first pass, unchecked against the pages by a human.
%
% Dating: \dating{[1970]} is the folder's own date in src/content/catalogue.ts,
% copied verbatim; since the folder has a date of its own, the group rule
% (« Descentes », 126-133) does not apply. Nothing on these leaves dates them.
%
% Pages skipped: 78, a blank cream sheet.
%
% Pages identified, not transcribed: all the others (61-77, 79-80). Nothing
% in this batch is a page written in hand throughout. Each page gets one
% short French \note identifying it, followed by every handwritten
% intervention on it, given in full; the typed text is never re-set.
%  - Pages 61-75: the photocopied seminar typescript continues (chapter III,
%    « Groupes de Barsotti-Tate », §3.2 to §7.4). Pencil folios « III-6 » ...
%    « III-20 », top right. Proof corrections on pages 62, 63, 65 and 69 (a
%    margin cross, « n+1 », « Sorites » for « Suites », « biunivoquement »
%    for « G univoquement », a hollow μ), of the same kind as batches 2-3.
%  - Page 76: a sheet of computer-listing paper carrying only a pencilled
%    label « I, II / Vecteurs de Witt et théorie de Dieudonné ».
%  - Page 77: a short typed list of chapter titles I-VI (Witt vectors and
%    Dieudonné theory; BT groups; divided powers and crystals; the Dieudonné
%    functor over an arbitrary base; F-crystals and a specialisation
%    theorem), renumbered by typed overstrikes, with two pencil insertions.
%  - Pages 79-80: an ORIGINAL typescript (ribbon, not a photocopy; another
%    machine than the seminar copies; many « xxxx » overstrikes), headed
%    « Rappels sur la théorie de Dieudonné (sur corps parfait k) », §1
%    « Rappels sur vecteurs de Witt », with many ink interventions (blue and
%    black): bibliography in the margin, « p ... premier fixé dans tout le
%    cours », hand-drawn arrows and labels in two commutative squares.
% Whose hand, and whose typing: not settled from the leaves, and the notes
% say so rather than attribute; hence no \marginal{} anywhere. FOR THE HUMAN
% CHECK: the typescript of 77 and 79-80 looks like a working draft (the
% overstrikes, « puiss. div. », « base ± quelconque », the hand-numbered
% chapter plan of the Montreal course) and may be Grothendieck's own typing.
% If a reader who knows his typescripts confirms it, pages 77 and 79-80
% should be transcribed in full in a revision; this pass applied the
% folder's rule (typed text of undetermined author is not re-set).
% Batch 5 (from page 81) presumably continues the same typescript.
%
% His pagination: none certain. Folios « III-6 »-« III-20 » in pencil,
% hand undetermined; « I et II » top right of page 79.
%
% What the batch is about. (1) Pages 61-75: Barsotti-Tate groups — 3.2 the
% definition of truncated BT groups of level n (condition (*) Ker F = Im V
% for n = 1), 3.3; §4 BT groups (p-divisible groups after Tate, co-p-adic
% systems, 4.1-4.2); §5 « Sorites » (5.1 description by co-p-adic systems,
% base change; 5.2 extensions and quotients, remark; 5.3 rank or height;
% 5.4 p-adic systems T_p(G); 5.5 the dual G*; 5.6 the Dieudonné module over
% a perfect field, D*(G) = lim M(n)); §6 examples (abelian schemes, ind-étale
% groups and twisted constant p-adic sheaves, tori and μ, extensions, toroidal
% groups, groups with connected fibres, formal Lie varieties and groups, the
% conjectures on formal Lie groups being BT); §7 the composition series
% (G^o, G^et, G^tor over a point; 7.3 failure over a general base, via the
% modular elliptic curve; 7.4 Proposition, conditions i)-iii)bis, and the
% Lemma on factorising a finite locally free morphism). (2) Pages 76-77: the
% label and the typed plan of a course in six chapters. (3) Pages 79-80:
% Witt vectors — W_n, W, the ghost components, V, R, T, F, FV = VF = p,
% the ind-scheme W→, the homomorphism u and W(k) for k perfect, u : W→(k) ≃ K/W.
%
% Hardest pages: 79 (the oblique marginal bibliography and the note
% « p ... premier fixé ») and 80 (hand-drawn arrows and labels in the two
% squares). \ill{}: 0. \uncertain{}: 3.

\documentclass[11pt,a4paper]{article}
% Le préambule commun à toutes les transcriptions du fonds Grothendieck.
%
% Il tient en une idée : **l'incertitude se compose, elle ne se résout pas.**
% Une transcription qui lisse un mot douteux a détruit la seule chose pour
% laquelle on la faisait — un lecteur doit pouvoir savoir, à chaque endroit,
% ce qui était sur le feuillet et ce qui a été ajouté. Les six macros qui
% suivent sont donc l'appareil critique, et non de la décoration.
%
% Les mêmes commandes sont reconnues par `scripts/render.mjs`, qui produit la
% vue de lecture du site. Une macro ajoutée ici doit l'être aussi là-bas,
% faute de quoi le rendu échouera bruyamment — ce qui est voulu : un rendu
% silencieusement faux est pire qu'un rendu absent.

\ProvidesPackage{grothendieck}[2026/08/08 Transcriptions du fonds Grothendieck]

\usepackage{amsmath,amssymb,amsthm}
\usepackage{mathrsfs}
\usepackage{stmaryrd}
% Les diagrammes commutatifs sont partout dans ce fonds : c'est la langue
% même de la géométrie algébrique de Grothendieck, et une transcription qui
% les rendrait en prose ne transcrirait rien.
\usepackage{tikz-cd}
% Français par défaut — c'est la langue des feuillets — mais l'anglais est
% chargé aussi : la traduction et le résumé sont des documents anglais, et la
% typographie française (espaces avant les deux-points) y serait une faute.
% Ces documents basculent par \selectlanguage{english} en tête de corps.
\usepackage[english,french]{babel}
\usepackage{fontspec}
\usepackage{geometry}
\geometry{a4paper,margin=25mm,marginparwidth=28mm}
\usepackage{marginnote}
\usepackage{xcolor}
% ulem, not soul. `soul`'s \st and \ul are built on a character-by-character
% scan that XeLaTeX's font machinery breaks: the document fails with an
% undefined control sequence rather than degrading. ulem does the same two jobs
% and is Unicode-engine safe. `normalem` stops it hijacking \emph, which the
% transcriptions use for Grothendieck's own emphasis.
\usepackage[normalem]{ulem}
\usepackage{hyperref}
\hypersetup{colorlinks=true,linkcolor=black,urlcolor=black,pdfborder={0 0 0}}

% --- Métadonnées du lot -----------------------------------------------------
% Reprises telles quelles par le script de rendu, qui les affiche en tête de
% la vue de lecture. Elles fixent ce que le fichier prétend couvrir : sans
% elles, une transcription flotte sans point d'attache dans un fonds de seize
% mille feuillets.

\newcommand{\folder}[1]{\gdef\grothFolder{#1}}
\newcommand{\batch}[1]{\gdef\grothBatch{#1}}
\newcommand{\leaves}[2]{\gdef\grothFirst{#1}\gdef\grothLast{#2}}
\newcommand{\foldertitle}[1]{\gdef\grothFoldertitle{#1}}
\gdef\grothFolder{?}\gdef\grothBatch{?}\gdef\grothFirst{?}\gdef\grothLast{?}\gdef\grothFoldertitle{}

% --- Le résumé --------------------------------------------------------------
%
% Il ouvre la lecture modernisée, sous le titre « Résumé », et il est composé
% en petit. Ce n'est pas un ornement : c'est le seul passage du document qui
% ne soit pas des mathématiques, et un lecteur doit pouvoir le reconnaître
% avant de l'avoir lu — puis le sauter s'il connaît déjà le sujet, ou s'y
% arrêter s'il ne le connaît pas. Le filet qui le suit marque la reprise du
% corps.

\newenvironment{resume}
  {\small\setlength{\parskip}{0.45em}}
  {\par\medskip\hrule\medskip}

% --- Mots-clés ---------------------------------------------------------------
%
% En anglais, en clôture du résumé de la lecture modernisée : le vocabulaire
% moderne sous lequel chercher ce que les feuillets construisent. Le manifeste
% du site (`npm run manifest`) les extrait pour étiqueter la cote entière —
% cette ligne est la source unique des tags, il n'y a pas de fichier à côté.
\newcommand{\keywords}[1]{\par\smallskip\noindent
  {\footnotesize\textsc{Keywords} --- \textit{#1}}\par}

% --- L'appareil critique ----------------------------------------------------

% Le numéro de feuillet, en marge. C'est l'ancre qui permet de régler un
% désaccord sur une formule en regardant *un* feuillet plutôt qu'une chemise
% de six cents. Le site s'en sert aussi pour tourner les pages du fac-similé
% quand on fait défiler la transcription.
\newcommand{\leaf}[1]{%
  \marginnote{\footnotesize\color{gray!70}\textsf{#1}}%
  \label{leaf:#1}\ignorespaces}

% Illisible : on ne devine pas. Un mot inventé qui se lit comme les autres est
% le pire résultat possible.
\newcommand{\ill}{\textcolor{red!60!black}{[\,\ldots\,]}}

% Lecture douteuse : le mot est donné, et signalé comme incertain.
\newcommand{\uncertain}[1]{\uline{#1}}

% Ajout de l'éditeur — une lettre manquante, un mot sous-entendu.
\newcommand{\add}[1]{\textcolor{blue!50!black}{[#1]}}

% Biffé par Grothendieck lui-même. Ce qu'il a rayé fait partie du document :
% une rature dit souvent où la pensée a changé de direction.
\newcommand{\struck}[1]{\sout{#1}}

% Note du transcripteur, sur l'état matériel du feuillet.
\newcommand{\note}[1]{\par\smallskip{\small\color{gray!60!black}\textit{[#1]}}\par\smallskip}

% Note marginale de Grothendieck, distincte du corps du texte.
\newcommand{\marginal}[1]{\par\smallskip\noindent\hspace{1em}%
  {\small\color{gray!60!black}#1}\par\smallskip}

% --- Titre ------------------------------------------------------------------

\AtBeginDocument{%
  \begin{center}
    {\large\bfseries Fonds Alexandre Grothendieck --- cote n\textsuperscript{o}~\grothFolder}\\[2pt]
    {\itshape\grothFoldertitle}\\[4pt]
    {\small feuillets \grothFirst--\grothLast\ (lot \grothBatch)}\\[2pt]
    {\footnotesize\color{gray!60!black}Transcription établie d'après le fac-similé numérisé
     par l'Université de Montpellier.\\
     Les lectures incertaines sont soulignées, les illisibles notées [\,\ldots\,],
     les ajouts entre crochets.}
  \end{center}
  \vspace{1em}}

\endinput


\folder{129}
\batch{4}
\pages{61}{80}
\dating{[1970]}
\watermark{Édition de démonstration}
\foldertitle{Notes séminaire Montréal (Delale, Labute, Hakim, Messing) : copies de tapuscrits annotés (s.d.), notes manuscrites (s.d.)}

\begin{document}

\page{61}
\note{tapuscrit, chapitre~III « Groupes de Barsotti-Tate », suite du §3
commencé page 60 : n°~3.2 Définition (groupe de Barsotti-Tate tronqué
d'échelon $n$ sur $S$, avec pour $n = 1$ la condition (*)
$\mathrm{Ker}\, F_{G_0/S_0} = \mathrm{Im}\, V_{G_0}$) et n°~3.3
($G(i)$, $1 \leq i \leq n$, est tronqué d'échelon $i$ ; diagramme
cartésien). Texte d'autrui, non transcrit. Seule intervention manuscrite,
d'une main qu'on ne peut pas identifier d'après le feuillet : le folio
«~III-6~», au crayon}

\page{62}
\note{même tapuscrit, §4 « Groupes de Barsotti-Tate » : terminologie
(renvoi à Tate, \emph{$p$-divisible groups}, Driebergen 1966), n°~4.1
Définition ($p$-divisible), systèmes co-$p$-adiques
$G_1 \hookrightarrow G_2 \hookrightarrow \dots$, et n°~4.2 Définition
(groupe de Barsotti-Tate ; renvoi à 2.4). Interventions manuscrites, main
non identifiée : folio «~III-7~», au crayon ; en marge une croix et
«~$n+1$~», devant le système inductif, dont l'indice $n+1$ de
$G_{n+1}$ est retouché à la main (le signe $+$ repassé)}

\page{63}
\note{même tapuscrit, fin du §4 (réciproque sur un corps algébriquement
clos) et §5, n°~5.1 (description par les systèmes co-$p$-adiques ;
changement de base) et n°~5.2 Proposition (extensions et quotients de
groupes de BT), début de la démonstration. Interventions manuscrites, main
non identifiée : folio «~III-8~», au crayon ; en marge une croix, et dans
le titre du §5 le mot « Suites » biffé d'un trait oblique, « Sorites » écrit
au-dessus, ce qui donne « Sorites sur les groupes de Barsotti-Tate »}

\page{64}
\note{même tapuscrit, fin de la démonstration de 5.2 (renvoi à II,~1.6),
Remarque (le noyau de $G \to G$ n'est pas en général un groupe de BT),
n°~5.3 « Rang (ou hauteur) d'un groupe de Barsotti-Tate »
($\mathrm{rang}\,(G(n)_s) = p^{n r(s)}$) et début du n°~5.4 (systèmes
$p$-adiques). Seule intervention manuscrite : le folio «~III-9~», au
crayon}

\page{65}
\note{même tapuscrit, suite du n°~5.4 ($T_p(G)$), n°~5.5 « Dual d'un
groupe de Barsotti-Tate » ($G^* = \varinjlim G(n)^*$, équivalence entre
$\mathrm{BT}(S)$ et son opposée) et début du n°~5.6 « Module de Dieudonné
associé à un groupe de BT sur un corps parfait » (renvois à II,~4.1 et II,
5.3.1). Interventions manuscrites, main non identifiée : folio «~III-10~»,
au crayon ; en marge une croix, et dans « correspondre G univoquement le
système projectif » les mots « G univoquement » biffés d'un trait oblique,
« biunivoquement » écrit au-dessus}

\page{66}
\note{même tapuscrit, suite du n°~5.6 : énoncé de l'équivalence $D^*$
entre $[\mathrm{BT}(k)]^{\circ}$ et les $\mathcal{D}$-modules libres de
type fini sur $W$, compatibilités à $F$ et $V$, modules
$M(n) = D^*(G(n))$ et définition $D^*(G) = \varprojlim M(n)$ (renvois à
II,~4 et 2.2). Seule intervention manuscrite : le folio «~III-11~», au
crayon}

\page{67}
\note{même tapuscrit, fin du n°~5.6 ($D^*(G)$ libre de type fini sur $W$,
muni de $F$ et $V$) et §6 « Exemples et types particuliers de groupes de
Barsotti-Tate », n°~6.1 (schémas abéliens ; renvois à Lang,
\emph{Abelian Varieties}, et à EGA~IV, 11.3.10 ;
$\mathrm{rang}\, A(1) = p^{2d}$). Seule intervention manuscrite : le folio
«~III-12~», au crayon}

\page{68}
\note{même tapuscrit, suite du n°~6.1 (cas ind-étale, dualité des variétés
abéliennes), n°~6.2 (groupes de BT ind-étales ; équivalence avec les
faisceaux $p$-adiques constants tordus sans torsion, monodromie de
$\pi_1(S, \bar s)$) et n°~6.3 Exemple (tores). Seule intervention
manuscrite : le folio «~III-13~», au crayon}

\page{69}
\note{même tapuscrit, suite du n°~6.3 ($T(\infty)$, groupe
$\mu = \mathbb{G}_{m,S}(\infty) = \varinjlim \mu_{p^n,S}$), n°~6.4 Exemple
(extension d'un schéma abélien par un tore, rang $2d+r$), n°~6.5 (groupes
de BT toroïdaux ; anti-équivalence avec les faisceaux $p$-adiques
constants tordus sans torsion). Interventions manuscrites, main non
identifiée : folio «~III-14~», au crayon ; en marge une croix et un $\mu$
ajouré, devant la formule qui définit $\mu$, où le $\mu$ de
$\varinjlim \mu_{p^n,S}$ est retracé à l'encre, ajouré lui aussi}

\page{70}
\note{même tapuscrit, fin du n°~6.5 (l'équivalence
$G \mapsto \underline{\mathrm{Hom}}_{\mathrm{gr}}(\mu, G)$), n°~6.6
(groupes de BT à fibres connexes ; renvoi à II,~2.2) et n°~6.7 (variété de
Lie formelle ponctuée sur $S$), Remarques, 1). Seule intervention
manuscrite : le folio «~III-15~», au crayon}

\page{71}
\note{même tapuscrit, fin des Remarques du n°~6.7, définition des groupes
de Lie formels, et conditions 1) ($G$ $p$-divisible) et 2) ($G(1)$
représentable par un schéma fini localement libre) pour qu'un groupe de Lie
formel soit de BT. Seule intervention manuscrite : le folio «~III-16~», au
crayon}

\page{72}
\note{même tapuscrit, conditions 3) et 4) sur la fibre géométrique dans le
cas artinien local, conjecture sur les $S$-morphismes de variétés de Lie
formelles (renvoi à SGA~3, VII, numéro laissé en blanc), exemple de la
courbe elliptique modulaire et de l'invariant de Hasse, et seconde
conjecture (rang des fibres localement constant). Seule intervention
manuscrite : le folio «~III-17~», au crayon}

\page{73}
\note{même tapuscrit, §7 « Suite de composition d'un groupe de
Barsotti-Tate » : n°~7.1 (base réduite à un point, $G^{\circ}$,
$G^{\mathrm{et}}$, suite exacte
$0 \to G^{\circ} \to G \to G^{\mathrm{et}} \to 0$) et n°~7.2 ($G^{\mathrm{tor}}$,
filtration $0 \subset G^{\mathrm{tor}} \subset G^{\circ} \subset G$). Seule
intervention manuscrite : le folio «~III-18~», au crayon}

\page{74}
\note{même tapuscrit, fin du n°~7.2 (isomorphismes
$(G/G^{\circ})^* \simeq (G^*)^{\mathrm{tor}}$,
$(G/G^{\mathrm{tor}})^* \simeq (G^*)^{\circ}$) et n°~7.3 (sur une base
quelconque, pas d'extension en général ; argument par le rang séparable,
exemple de la variété elliptique modulaire). Seule intervention
manuscrite : le folio «~III-19~», au crayon}

\page{75}
\note{même tapuscrit, fin du n°~7.3 et n°~7.4 Proposition (équivalence des
conditions i), ii), ii)$^{\mathrm{bis}}$, iii), iii)$^{\mathrm{bis}}$), début
de la démonstration et Lemme (factorisation d'un morphisme fini localement
libre en un morphisme radiciel et surjectif suivi d'un morphisme étale ;
renvoi à EGA~IV, numéro laissé en blanc). Seule intervention manuscrite :
le folio «~III-20~», au crayon}

\page{76}
\note{feuillet de papier listing, sans texte dactylographié ni
mathématiques ; seule une inscription au crayon, en haut à droite, d'une
main qu'on ne peut pas identifier d'après le feuillet : «~I, II~» puis
«~Vecteurs de Witt et théorie de Dieudonné~». Elle semble désigner le
texte qui suit (la page 79 porte «~I et II~» en tête)}

\page{77}
\note{feuillet dactylographié original (et non une photocopie), d'une
autre machine que les tapuscrits précédents : une liste de titres de
chapitres numérotés I à VI (rappels sur les vecteurs de Witt et la
théorie de Dieudonné ; généralités sur les groupes de BT ; puissances
divisées et cristaux ; foncteur de Dieudonné sur une base quelconque,
« programme préliminaire » ; F-cristaux et un théorème de spécialisation
pour les groupes de BT). La numérotation a été décalée à la machine, par
surfrappe des chiffres romains. Auteur non indiqué ; texte non transcrit.
Interventions manuscrites, au crayon, main non identifiée : un «~II~»
inséré au-dessus de la première ligne, entre « et » et « la théorie de
Dieudonné », avec une parenthèse qui sépare les deux sujets (ce qui fait
de la théorie de Dieudonné le chapitre~II) ; au bout de la ligne~V, après
« programme préliminaire », l'ajout «~et problèmes~»}

\page{79}
\note{tapuscrit original (ruban, nombreuses surfrappes « xxxx »), de la
même machine que la page 77, intitulé « Rappels sur la théorie de
Dieudonné (sur corps parfait $k$) », §1 « Rappels sur vecteurs de Witt » :
schémas en anneaux $\underline{W}_n$, $\underline{W} = \varprojlim
\underline{W}_n$, composantes fantômes $\phi_i$, morphismes $V$, $R$, $T$
(avec $R_nT_n = V_n$, $T_nR_n = V_{n+1}$) et $F$, cas de la
caractéristique $p$. Auteur non indiqué ; texte non transcrit.
Interventions manuscrites, à l'encre noire et bleue, main non identifiée :
en haut à droite «~I et II~» ; en marge gauche, en oblique et soulignés,
deux renvois : «~Serre, Corps Locaux~» et
«~\uncertain{Demazure-Gabriel}, Schémas en Groupes~» ; à droite du titre
du §1, relié à lui par un long trait, «~$p$ \uncertain{nb} premier fixé
dans tout le cours~» ; dans « $(n \geq 1)$ » le signe $\geq$ écrit à
l'encre bleue, la parenthèse prolongée par un trait vers
$\underline{W}_n$ ; devant « $: \underline{W}_n \to \underline{O}^n$ »,
l'ajout «~$\Phi = (\phi_1, \dots, \phi_n)$~», et les exposants $n$ de
$\underline{O}^n$ et de $\underline{E}^n$ écrits ou repassés à la main ;
après le titre souligné « Morphismes remarquables : », l'ajout «~Morphismes
\emph{additifs}~» ; « puis » suivi d'un mot manuscrit biffé,
\struck{morphi} ; devant la ligne de $R_n$, «~morphisme d'Anneaux~», et
devant celle de $T_n$, «~morphisme add.~» ; quelques virgules ajoutées en
bout de ligne}

\page{80}
\note{suite du même tapuscrit original : $F$ homomorphisme d'anneaux
(Frobenius habituel), $V$ Verschiebung, $FV = VF = p \cdot \mathrm{id}$ en
caractéristique $p$ ; le schéma en groupes $\underline{W}_{\to} =
\varinjlim \underline{W}_n$ (morphismes de transition du type $D$),
$\underline{W}_{\to} = \underline{E}^{(\mathbb{N})}$ ; deux carrés
commutatifs en caractéristique $p$ ; l'homomorphisme canonique
$u : \underline{W}_{\to} \to \varinjlim (\underline{W}_n, p)$ ; pour $k$
parfait, $W = \underline{W}(k)$ anneau de valuation discrète, $K = W[1/p]$,
et l'isomorphisme $u : \underline{W}_{\to}(k) \simeq K/W$ ; la page s'arrête
en cours de phrase. Texte non transcrit. Interventions manuscrites, à
l'encre noire et bleue, main non identifiée : les flèches soulignant les
$\underline{W}_{\to}$ et la petite flèche sous le premier « lim » sont
tracées à la main ; dans le premier carré ($\underline{W}_n \to
\underline{W}_{n+1}$ par $T$), les deux flèches verticales montantes sont
tracées à la main, celle de droite étiquetée «~id~», et une lettre
manuscrite, lue \uncertain{r}, est posée sur la flèche du bas ; dans le
second carré, les deux flèches verticales descendantes sont tracées à
l'encre bleue, étiquetées $F^{n-1}$ et $F^{n}$ (l'exposant de droite
empâté) ; dans la définition de $u$, un indice $\mathbb{F}_p$ ajouté sous
$\underline{W}_{\to}$ et sous $\underline{W}_n$, ce qui donne
$u : \underline{W}_{\to\,\mathbb{F}_p} \to \varinjlim
(\underline{W}_{n\,\mathbb{F}_p}, p)$ ; la flèche $\simeq$ de
$u : \underline{W}_{\to}(k) \simeq K/W$ est tracée à la main}

\end{document}
